% NOTE: PB1/CM numbers are from the pilot (100 files, one seed; 6.1's prediction 500 files), update after the scale-up;
% AU's evasion numbers are final for seed 0 (runs 47-49: 902 detected held-out malware, loader-safe changes).
Static malware detectors come in two kinds: models over engineered features, such as gradient-boosted trees on EMBER
features, and deep networks that read a file's raw bytes. Their explanations live in different spaces. A feature-based
detector's explanation names features without saying where in the file their evidence lies; a raw-byte detector's
explanation marks bytes that have no name and mean something different in every file. We observe that every engineered
feature is computed from specific bytes, and record this by instrumenting the feature extractor as a \emph{provenance
matrix}: each byte's share of each feature, grounded in the Shapley value. This one matrix transfers explanations in
both directions. The \emph{pullback} places feature attributions on the bytes, exactly for linear models over additive
features; the \emph{pushforward} places byte attributions in named features, a first step toward inverting the
pullback that already comes close to the exact inverse. On EMBER-2024 with a gradient-boosted tree and MalConv, the
pullback is as faithful as window occlusion at about a hundredth of its time and, at byte resolution, removes more
evidence than any direct explanation, from one feature extraction instead of tens of thousands. Comparing the two
detectors file by file, it predicts which of the tree's detections an attacker guided only by MalConv also weakens
(Spearman 0.48, against 0.21 for occlusion). Pushed forward, MalConv's per-file explanations become named, recurring
reliance, about twice as faithful as a surrogate model's: changing the confirmed fields only as the Windows loader
allows evades 12.4\% of detected held-out malware, twice the rate of other loader-ignored fields, and zeroing them
before scoring halves the strongest such attack.
